\documentclass[12pt]{article}

% ---------------- Packages ----------------
\usepackage[margin=1in]{geometry}
\usepackage{amsmath, amssymb}
\usepackage{graphicx}
\usepackage[margin=1in]{geometry}
\usepackage{float}
\usepackage{booktabs}
\usepackage{caption}
\usepackage{subcaption}
\usepackage{tabularx}
\usepackage{siunitx}
\usepackage{setspace}
\usepackage{hyperref}

\setlength{\parindent}{0pt}
\onehalfspacing
\begin{document}

\begin{titlepage}
    \centering

    \vspace*{1in}

    \textbf{\Large COLLEGE OF LAKE COUNTY} \\[0.2in]

    \textbf{\large Department of Engineering} \\[0.6in]

    \textbf{\Large EGR 260} \\
    \textbf{\large Introduction to Circuit Analysis} \\[0.6in]

    \textbf{\Large Lab \#1} \\[0.2in]


    \begin{center}
            \textbf{Reported by:} Anh Thong Le \\[0.15in]
    \textbf{Date Performed:} January 29, 2026 \\[0.15in]
    \textbf{Date Submitted:} February 11, 2026 \\[0.15in]
    \textbf{Submitted to:} Dr. Regez
    \end{center}
        

    \vfill
\end{titlepage}


% ---------------- Title Info ----------------
\title{\textbf{EGR 260: Intro to Circuit Analysis}\\
Lab \#1: Ohm's Law and Basic Circuit Analysis}
\author{
Your Name \\
Instructor: Instructor Name \\
Department of Engineering \\
University Name
}
\date{Date Performed: \quad Date Submitted:}

% ---------------- Objective ----------------
\section{Objective}
The objective of this experiment is to verify Ohm’s Law by analyzing the relationship between voltage, current, and resistance in a simple electrical circuit. Experimental results are compared with theoretical and simulated values to evaluate measurement accuracy and sources of error.

% ---------------- Equipment ----------------
\section{Equipment}
\begin{itemize}
    \item DC power supply
    \item Digital multimeter
    \item Resistors (specified values)
    \item Breadboard
    \item Connecting wires
\end{itemize}

% ---------------- Experimental Procedure ----------------
\section{Experimental Procedure}

\subsection{Circuit Schematic}
Figure~\ref{fig:schematic} shows the schematic of the circuit used in this experiment.

\begin{figure}[H]
    \centering
    \includegraphics[width=0.5\textwidth]{image.png}
    \caption{Schematic of simple resistive circuit}
    \label{fig:schematic}
\end{figure}

\subsection{Procedure}
\begin{enumerate}
    \item The circuit was constructed according to the schematic shown in Figure~\ref{fig:schematic}.
    \item A DC voltage was applied using the power supply.
    \item Voltage across the resistor and current through the circuit were measured using a multimeter.
    \item Measurements were repeated for multiple voltage levels.
\end{enumerate}

% ---------------- Analysis ----------------
\section{Analysis}

\subsection{Sample Ohm’s Law Calculation}
Ohm’s Law is given by:
\begin{equation}
    V = IR
\end{equation}

Given:
\[
V = \SI{1}{\volt}, \quad R = \SI{1000}{\ohm}
\]

The theoretical current is:
\begin{equation}
    I = \frac{V}{R} = \frac{1}{1000} =  \SI{0.001}{\ampere}
\end{equation}

% ---------------- Data ----------------
\section{Data}
\subsection{Measurement of resistor}
\begin{itemize}
    \item R1 = 0.988k\SI{}{\ohm} normally 1k\SI{}{\ohm}
    \item R2 = 2.155k\SI{}{\ohm} normally 2k\SI{}{\ohm}

    \item R3 = 2.856k\SI{}{\ohm} normally 3k\SI{}{\ohm}
    \item R4 = 9.840k\SI{}{\ohm} normally 10k\SI{}{\ohm}
    
\end{itemize}

\subsection{Simulation and measurement values}
\begin{table}[H]
\centering
\begin{tabularx}{0.8\textwidth} { 
  | >{\centering\arraybackslash}X 
  | >{\centering\arraybackslash}X 
  | >{\centering\arraybackslash}X | }
 \hline
 Voltage (V) & Current (A) in the R1 circuit simulated & Current (A) in the R1 circuit measured \\
 \hline
 0   & 0   & 0   \\
\hline
 1   & 0.001 & 1.002$\times 10^{-3}$  \\
\hline
 2   & 0.002 & 2.034$\times 10^{-3}$   \\
\hline
 3   & 0.003 & 3.002$\times 10^{-3}$   \\
\hline
 4  & 0.004  & 4.075$\times 10^{-3}$  \\
\hline
 5   & 0.005  & 5.042$\times 10^{-3}$  \\
\hline
 6 & 0.006 & 6.050$\times 10^{-3}$   \\
\hline
 7 & 0.007  & 7.066$\times 10^{-3}$   \\
\hline
 8 & 0.008 & 8.091$\times 10^{-3}$   \\
\hline
 9 & 0.009  & 9.161$\times 10^{-3}$   \\
\hline
 10 & 0.01 & 10.129$\times 10^{-3}$   \\
\hline
\end{tabularx} 
\caption{Voltage and current values for simple 1k\SI{}{\ohm} resistor circuit}

\end{table}

\begin{table}[H]
\centering
\begin{tabularx}{1\textwidth} { 
  | >{\centering\arraybackslash}X 
  | >{\centering\arraybackslash}X 
  | >{\centering\arraybackslash}X |
   >{\centering\arraybackslash}X |
    >{\centering\arraybackslash}X |
     >{\centering\arraybackslash}X |
      >{\centering\arraybackslash}X |}
 \hline
 Voltage (V) & Current (A) in the R2 circuit simulated & Current (A) in the R2 circuit measured & Current (A) in the R3 circuit simulated & Current (A) in the R3 circuit measured & Current (A) in the R4 circuit simulated & Current (A) in the R4 circuit measured \\
 \hline
 0 & 0  & 0& 0  & 0  & 0  & 0  \\
\hline
 2   & 0.001  & 0.954$\times 10^{-3}$  & 0.667$\times 10^{-3}$ & 0.668$\times 10^{-3}$  & 0.2$\times 10^{-3}$ &0.208$\times 10^{-3}$\\
\hline 
 4   & 0.002  & 1.854$\times 10^{-3}$  & 1.333$\times 10^{-3}$ &1.348$\times 10^{-3}$ & 0.4$\times 10^{-3}$ &0.414$\times 10^{-3}$ \\
\hline
 6   & 0.003  & 2.819$\times 10^{-3}$  & 2.000$\times 10^{-3}$ & 2.052$\times 10^{-3}$  &0.6$\times 10^{-3}$  & 0.612$\times 10^{-3}$  \\
\hline
 8 & 0.004  & 3.737$\times 10^{-3}$  &2.667$\times 10^{-3}$ & 2.72$\times 10^{-3}$ & 0.8$\times 10^{-3}$& 0.822$\times 10^{-3}$ \\
\hline
 10   & 0.005  & 4.687$\times 10^{-3}$ &3.333$\times 10^{-3}$ & 3.41$\times 10^{-3}$&1.000$\times 10^{-3}$ & 1.030$\times 10^{-3}$ \\
\hline
\end{tabularx} 
\caption{Voltage and current values for simple 2k\SI{}{\ohm},  3k\SI{}{\ohm},  10k\SI{}{\ohm} resistor circuit}

\end{table}

% ---------------- Graphs ----------------
\section{Graphs}

% \begin{figure}[H]
%     \centering
%     \includegraphics[width=0.7\textwidth]{iv_graph.png}
%     \caption{Current vs. Voltage relationship}
% \end{figure}

\begin{figure}[H]
    \centering
    \includegraphics[width=0.7\linewidth]{linearRelationshipR1Edited.png}
    \caption{Voltage vs. measured current graph for R1}
    \label{fig:placeholder}
\end{figure}

\begin{figure}[H]
    \centering
    \includegraphics[width=0.7\linewidth]{linearRelationshipR234.png}
    \caption{Voltage vs. measured current graph for R2, R3, and R4}
    \label{fig:placeholder2}
\end{figure}

% ---------------- Discussion ----------------
\section{Discussion and Conclusions}

\subsection{Relationship}
The experimental results in figure \ref{fig:placeholder} and figure \ref{fig:placeholder2} demonstrate a linear relationship between voltage and current, confirming that current is linearly proportional to voltage and inversely proportional to resistance. The slope of the voltage-measured current graph represents the values of resistance.

\subsection{Experimental and Simulated Results}
The measured values were close to simulated results, with minor discrepancies.

Possible sources of error include:
\begin{itemize}
        \item The difference in initial values of resistors 
    \item Internal resistance of the multimeter
    \item Contact resistance from wiring and breadboard connections
\end{itemize}

\subsection{Conclusion}
This experiment successfully verified Ohm’s Law through experimental measurement. The results closely matched theoretical predictions within small discrepancies.

\end{document}
